% This file is part of the modular source for Stratified Relational Workspace Theory.
% !TeX root = ../Stratified_Relational_Workspace_Theory.tex
\theorypart{XXII}{Integrated Worked Examples and Case Studies}

\topic{77}{Case-Study Method and Common Gate Card}

The preceding parts define a large family of structural, operational, causal,
relational, persistent, formative, and agentive predicates. A worked example
is useful only when it makes clear which of those predicates are established,
which are assumed, and which remain unmeasured. This part therefore uses a
single case-study protocol for three different kinds of realization:
\begin{enumerate}
\item an exact finite-dimensional model of a human--human theorem laboratory;
\item an idealized human--AI--artifact--executor research loop, followed by a
 deliberately conservative audit of the collaboration that motivated the
 theory; and
\item a closed-form multi-agent control system that reaches the formal agency
 level inside a small declared domain.
\end{enumerate}
The first and third cases are analytic models. Their conclusions hold inside
those models. The second contains both an analytic idealization and an
observational audit. The audit is not upgraded to a certificate merely because
the idealization is mathematically clean.

\begin{definition}[Case-study realization]\label{def:case-study-realization}
A case-study realization is a tuple
\[
\mathfrak K
=
\bigl(
\mathfrak R,
\mathfrak B,
\mathfrak C,
\mathfrak P,
\mathcal E,
\mathfrak D,
\mathfrak S,
\mathcal I,
\mathcal Q,
\mathcal G
\bigr),
\]
where:
\begin{itemize}
\item \(\mathfrak R\) is the declared relational realization;
\item \(\mathfrak B\) is the boundary problem and component decomposition;
\item \(\mathfrak C\) is the structural core--transfer model;
\item \(\mathfrak P\) is the downstream predictive or control problem;
\item \(\mathcal E\) is the evidence or generated-data family;
\item \(\mathfrak D\) is the admissible decomposition family;
\item \(\mathfrak S\) is the admissible family of complete scientific
 specifications, including targets, losses, horizons, interventions, nulls,
 and thresholds;
\item \(\mathcal I\) is the feasible intervention family;
\item \(\mathcal Q\) is the matched null family; and
\item \(\mathcal G\) is the set of gates being claimed.
\end{itemize}
A conclusion is always relative to this tuple. Changing the downstream target,
temporal horizon, intervention support, component decomposition, null family,
or threshold rule creates a different specification of the case. A stronger
case study reports whether the conclusion survives the predeclared family
\(\mathfrak S\) and whether the relevant gates are identified by the available
experiment family.
\end{definition}

\subtopic{77.1 Three evidential statuses}

Every claim in this part receives one of three primary statuses:
\begin{description}
\item[Analytic.] The gate follows exactly from a declared mathematical model.
\item[Empirically certified.] The gate is estimated from held-out data using
 the protocol of Part~XXI and passes its frozen uncertainty rule.
\item[Observational only.] The case contains suggestive records or natural
 perturbations, but no valid held-out intervention-and-null certificate.
\end{description}
An analytic result is not automatically an empirical result. An observational
pattern is not automatically a causal result. None of the three statuses
establishes phenomenal consciousness.

\subtopic{77.2 Common gate card}

For each case we report the vector
\[
\boxed{
\mathsf{Card}(\mathfrak K)
=
\left(
\mathsf S,
\mathsf O,
\mathsf R,
\mathsf P,
\mathsf F,
\mathsf A,
\mathsf D,
\mathsf X,
\mathsf C
\right),
}
\]
where:
\[
\begin{array}{ll}
\mathsf S & = \text{structural core and spectral coherence},\\
\mathsf O & = \text{operational alignment and causal efficacy},\\
\mathsf R & = \text{relational irreducibility and matched-null gap},\\
\mathsf P & = \text{persistence, reconstruction, and identity},\\
\mathsf F & = \text{formation provenance and boundary status},\\
\mathsf A & = \text{goal-conditioned agency and self-maintenance},\\
\mathsf D & = \text{decomposition robustness and interface audit},\\
\mathsf X & = \text{specification robustness, identification, and uniform uncertainty},\\
\mathsf C & = \text{consciousness or personhood claim}.
\end{array}
\]
The decomposition coordinate distinguishes a result established only for the
declared interface from one stable across an admissible family. The
specification coordinate records whether the same conclusion survives
scientifically defensible alternatives, whether its sign is identified by the
available experiments, and whether its lower margin remains positive after
uniform finite-sample uncertainty. The final coordinate is included to make
the non-implication explicit. In all examples below,
\[
\mathsf C=\text{not established}.
\]

\subtopic{77.3 Claim-level discipline}

Let \(L(\mathfrak K)\in\{0,1,\ldots,6\}\) denote the highest level from
Part~XXIV whose complete required gate set is established. The level is
conjunctive and pointwise in the declared scientific specification unless the
\(\mathsf D\) and \(\mathsf X\) coordinates also pass:
\[
L(\mathfrak K)\geq \ell
\quad\Longleftrightarrow\quad
\bigwedge_{j\in\mathcal J(\ell)}\Pass_j(\mathfrak K).
\]
A case may have strong higher-level-looking behavior while remaining at a lower
certified level because an earlier gate is absent. For example, successful
repair behavior does not establish self-maintenance when no matched no-repair
counterfactual is available, and coherent shared documents do not establish a
relational workspace when one participant contains the complete operative
state.

\subtopic{77.4 Structural calculations used repeatedly}

The positive examples use a recurring distributed code. Let
\[
c=(c_1,c_2,c_3)^{\mathsf T}\in\mathbb R^3
\]
and define three local encodings
\[
E_1c=(c_1,c_2),
\qquad
E_2c=(c_2,c_3),
\qquad
E_3c=(c_1,c_3).
\]
Then
\[
\sum_{i=1}^3E_i^*E_i=2I_3.
\]
No individual encoding is injective, while every two-node coalition is
injective. The pairwise frame operators are
\[
\begin{aligned}
E_{\{1,2\}}^*E_{\{1,2\}}&=\operatorname{diag}(1,2,1),\\
E_{\{1,3\}}^*E_{\{1,3\}}&=\operatorname{diag}(2,1,1),\\
E_{\{2,3\}}^*E_{\{2,3\}}&=\operatorname{diag}(1,1,2).
\end{aligned}
\]
Hence
\[
\alpha_{\{i,j\}}=1
\qquad(i\neq j),
\]
and one-node erasure recovery is exact. Every nonzero core vector is visible
at at least two nodes, so
\[
d_{\mathrm{er}}=2.
\]
This code is distributed but redundant: no node is sufficient, yet any one
node may be lost without destroying the core.

\subtopic{77.5 Binary synergy test}

Let \(C_1,C_2,C_3\) be independent uniform bits and let
\[
Y=C_1\oplus C_2\oplus C_3.
\]
Each node sees two of the three bits. Therefore
\[
H(Y\mid E_iC)=1\text{ bit}
\qquad\text{for every }i,
\]
while
\[
H(Y\mid E_1C,E_2C,E_3C)=0.
\]
Under logarithmic loss,
\[
\Risk_{\log}(\{i\})=1,
\qquad
\Risk_{\log}(V)=0,
\qquad
\Delta_{\mathrm{joint}}^{\log}=1\text{ bit},
\qquad
k_0=2.
\]
Moreover, \(Y\) is independent of every individual local state. The null law
that keeps all three core bits independent and makes \(Y\) an independent
uniform bit preserves every pair \((Y,E_iC)\) while destroying the joint parity
relation. It is therefore a predictive-marginal-matched relational null with
\[
\Delta_{\mathrm{tr}}=1\text{ bit}.
\]
This test will be reused with different semantic interpretations.

\topic{78}{Case Study A: A Distributed Human--Human Theorem Laboratory}

\subtopic{78.1 Declared realization}

Consider two researchers and one persistent notebook:
\[
V_{\mathrm{HH}}=\{P,C,N\},
\]
where:
\begin{itemize}
\item \(P\) is a proposer who tracks the theorem statement and open proof
 obligations;
\item \(C\) is a critic who tracks proof obligations and falsification status;
\item \(N\) is a notebook that tracks the theorem statement and falsification
 status.
\end{itemize}
The abstract core is
\[
c=(q,p,f)^{\mathsf T}\in\mathbb R^3,
\]
with
\[
\begin{aligned}
q&=\text{statement coordinate},
&
p&=\text{proof-obligation coordinate},\\
f&=\text{falsification-status coordinate}.
\end{aligned}
\]
The accessible encodings are
\[
E_P=
\begin{pmatrix}
1&0&0\\
0&1&0
\end{pmatrix},
\quad
E_C=
\begin{pmatrix}
0&1&0\\
0&0&1
\end{pmatrix},
\quad
E_N=
\begin{pmatrix}
1&0&0\\
0&0&1
\end{pmatrix}.
\]
The semantics are stylized. The point is not that real researchers contain
only two numbers. The point is to exhibit, without metaphor, a system in which
no participant or artifact contains the entire operative state while the
coupled system does.

\subtopic{78.2 Exact overlap descent}

Order the local coordinates as
\[
x=(q_P,p_P,p_C,f_C,q_N,f_N)^{\mathsf T}.
\]
Compatibility requires
\[
p_P=p_C,
\qquad
f_C=f_N,
\qquad
q_N=q_P.
\]
The discrepancy operator is
\[
D_{\mathrm{HH}}
=
\begin{pmatrix}
0&1&-1&0&0&0\\
0&0&0&1&0&-1\\
-1&0&0&0&1&0
\end{pmatrix}.
\]
Therefore
\[
\ker D_{\mathrm{HH}}
=
\left\{
(q,p,p,f,q,f)^{\mathsf T}:q,p,f\in\mathbb R
\right\}
\cong\mathbb R^3.
\]
The rows are orthogonal with squared norm two, so
\[
D_{\mathrm{HH}}D_{\mathrm{HH}}^*=2I_3.
\]
Consequently the transfer Laplacian has spectrum
\[
\boxed{
\spec(L_{\mathrm{HH}})
=
\{0,0,0,2,2,2\},
}
\]
and the exact workspace gap is
\[
\gamma_{\mathrm{HH}}=2.
\]
The low-energy subspace is therefore stable under any Laplacian perturbation
whose operator norm is strictly below one, with the quantitative projector
bound from Theorem~\ref{thm:spectral-stability}.

\subtopic{78.3 Reconstruction and one-node failure}

The full frame operator is
\[
F_{PCN}=E_P^*E_P+E_C^*E_C+E_N^*E_N=2I_3.
\]
For every two-node coalition \(S\),
\[
\lambda_{\min}(F_S)=1.
\]
Thus every two-node coalition reconstructs the complete core and, for noisy
observations \(y_S=A_Sc+\eta_S\),
\[
\|\widehat c-c\|
\leq
\|\eta_S\|.
\]
No individual node reconstructs \(c\), because
\[
\operatorname{rank}(E_P)
=
\operatorname{rank}(E_C)
=
\operatorname{rank}(E_N)
=2<3.
\]
The realization is therefore simultaneously:
\[
\text{individually incomplete}
\quad\text{and}\quad
\text{robust to any one-node loss}.
\]

\subtopic{78.4 Relational prediction problem}

Let \(q,p,f\) now be independent uniform bits and define the downstream
acceptance decision
\[
Y=q\oplus p\oplus f.
\]
The interpretation is abstract: acceptance depends on the statement, the state
of the proof obligations, and the current falsification record. No individual
node has sufficient information:
\[
H(Y\mid X_P)
=
H(Y\mid X_C)
=
H(Y\mid X_N)
=1.
\]
Every two-node coalition has all three bits and therefore zero conditional
entropy. Hence
\[
\boxed{
\Delta_{\mathrm{joint}}^{\log}=1\text{ bit},
\qquad
k_0=2.
}
\]
Let \(Q_{\mathrm{HH}}\) preserve the independent uniform local bit pairs but
replace \(Y\) by an independent uniform bit. Because parity is independent of
every two-bit local view, \(Q_{\mathrm{HH}}\) preserves each
\((Y,X_i)\)-marginal. Nevertheless,
\[
\Risk_{\log,Q_{\mathrm{HH}}}(V)
-
\Risk_{\log,P_{\mathrm{HH}}}(V)
=1\text{ bit}.
\]
This is an exact relational witness: the predictive capacity lies in the joint
organization rather than in any local predictive marginal.

\subtopic{78.5 Feasible intervention and specificity}

A feasible statement intervention must alter both local realizations of the
statement coordinate:
\[
\iota_{q\leftarrow a}:
(q_P,q_N)\leftarrow(a,a),
\]
while holding \((p,f)\) fixed. For configurations with fixed \(p\) and \(f\),
changing \(a\) flips the downstream decision with total-variation effect one:
\[
\TV
\left(
P(Y\mid\iota_{q\leftarrow0}),
P(Y\mid\iota_{q\leftarrow1})
\right)
=1.
\]
A matched intervention on a residual notebook-format coordinate leaves \(Y\)
unchanged. Thus the core intervention has positive causal specificity in this
model.

\subtopic{78.6 Persistence, handoff, and replacement}

Suppose the proposer becomes unavailable. The remaining coalition \(\{C,N\}\)
reconstructs \((q,p,f)\) exactly. A replacement proposer \(P'\) can therefore
be initialized by
\[
E_{P'}c=(q,p).
\]
If \(P'\) uses a different local coordinate system
\[
J_{P'}\in GL(2),
\]
then the incoming and outgoing overlap maps are conjugated accordingly. The
global section and all downstream kernels remain unchanged. This is exact
replacement invariance, not literal carrier identity.

If the notebook alone survives, the core cannot be reconstructed because the
proof-obligation coordinate is missing. The model therefore predicts a sharp
handoff boundary:
\[
\text{one surviving node is insufficient},
\qquad
\text{any two surviving nodes are sufficient}.
\]

\subtopic{78.7 Case card}

Inside the declared finite model,
\[
\boxed{
\begin{aligned}
\mathsf S&=\checkmark_{\mathrm{analytic}},
&\mathsf O&=\checkmark_{\mathrm{analytic}},\\
\mathsf R&=\checkmark_{\mathrm{analytic}},
&\mathsf P&=\checkmark_{\mathrm{analytic}},\\
\mathsf F&=\text{scaffolded},
&\mathsf A&=\text{not tested},\\
\mathsf D&=\text{fixed decomposition},
&\mathsf C&=\text{not established}.
\end{aligned}
}
\]
The highest supported level is Level~4 when the handoff dynamics and identity
correspondence are included. The example does not contain an operational goal
state, counterfactual policy, damage monitor, or repair loop, so it does not
support Levels~5 or 6.

\topic{79}{Case Study B: A Human--AI--Artifact--Executor Research Loop}

\subtopic{79.1 Idealized four-node loop}

Consider four heterogeneous components:
\[
V_{\mathrm{HAI}}=\{H,M,X,A\},
\]
where:
\begin{itemize}
\item \(H\) is a human researcher who maintains goals and semantic direction;
\item \(M\) is a language-model process that constructs local formalizations;
\item \(X\) is an execution and verification process that compiles, tests, or
 falsifies candidate artifacts; and
\item \(A\) is a persistent artifact store containing exact source, results,
 and versioned handoff state.
\end{itemize}
Let the abstract research state be
\[
c=(g,d,p,v)^{\mathsf T}\in\mathbb R^4,
\]
where:
\[
\begin{array}{ll}
g=\text{goal and acceptance criterion},
&d=\text{current derivation or theory state},\\
p=\text{open proof or implementation obligations},
&v=\text{verification and regression status}.
\end{array}
\]
Define local encodings
\[
E_Hc=(g,d),
\qquad
E_Mc=(g,p),
\qquad
E_Xc=(p,v),
\qquad
E_Ac=(d,v).
\]
Each coordinate is represented at exactly two nodes, and the incidence pattern
forms a four-cycle.

\subtopic{79.2 Exact structural calculation}

Order the local coordinates as
\[
x=(g_H,d_H,g_M,p_M,p_X,v_X,d_A,v_A)^{\mathsf T}.
\]
The overlap discrepancies are
\[
g_H-g_M,
\qquad
p_M-p_X,
\qquad
v_X-v_A,
\qquad
d_A-d_H.
\]
Thus
\[
D_{\mathrm{HAI}}
=
\begin{pmatrix}
1&0&-1&0&0&0&0&0\\
0&0&0&1&-1&0&0&0\\
0&0&0&0&0&1&0&-1\\
0&-1&0&0&0&0&1&0
\end{pmatrix}.
\]
The rows are orthogonal with squared norm two. Therefore
\[
\ker D_{\mathrm{HAI}}
\cong\mathbb R^4,
\]
\[
\boxed{
\spec(L_{\mathrm{HAI}})
=
\{0,0,0,0,2,2,2,2\},
\qquad
\gamma_{\mathrm{HAI}}=2.
}
\]
The full accessible frame operator is
\[
F_{HMXA}=2I_4.
\]
No single node is injective. Every three-node coalition reconstructs the full
state, and the two opposite pairs
\[
\{H,X\},
\qquad
\{M,A\}
\]
are already sufficient. Adjacent pairs observe only three of the four
coordinates. Hence
\[
k_{\mathrm{rec}}=2,
\qquad
d_{\mathrm{er}}=2,
\]
while the identity of the sufficient two-node coalition depends on the graph
geometry.

\subtopic{79.3 Relational parity benchmark}

Let \(g,d,p,v\) be independent uniform bits and define
\[
Y=g\oplus d\oplus p\oplus v.
\]
Every single node sees only two bits, so
\[
H(Y\mid X_i)=1
\qquad\text{for all }i.
\]
The opposite pairs \(\{H,X\}\) and \(\{M,A\}\) determine all four bits and
therefore determine \(Y\). Every adjacent pair remains one bit short. It
follows that
\[
\boxed{
\Delta_{\mathrm{joint}}^{\log}=1\text{ bit},
\qquad
k_0=2.
}
\]
As in the preceding case, parity is independent of every individual local
view. Replacing \(Y\) by an independent uniform bit produces a matched local
predictive null and a full-system transfer gap of one bit.

This benchmark is intentionally narrow. It demonstrates that the four-cycle
can realize exact relational information. It does not show that every real
human--AI collaboration has this encoding, this target, or this degree of
irreducibility.

\subtopic{79.4 A process-level interpretation}

The idealized loop can be read as the recurrent sequence
\[
\boxed{
H\xrightarrow{g}M
\xrightarrow{p}X
\xrightarrow{v}A
\xrightarrow{d}H.
}
\]
A more explicit cycle is:
\begin{enumerate}
\item the human supplies a goal and acceptance criterion;
\item the model converts that goal into a formal derivation and explicit
 obligations;
\item the executor compiles, tests, or searches for counterexamples;
\item the artifact store preserves the exact source and verification record;
\item the human evaluates the returned derivation and evidence, revising the
 next goal.
\end{enumerate}
The operative state is reentrant because products of one cycle become inputs to
the next. The artifact is not merely a record when its contents alter the next
formalization, test, or acceptance decision.

\subtopic{79.5 Conservative audit of the motivating collaboration}

The collaboration that motivated this theory has several features consistent
with the four-node description: specialized shared vocabulary, versioned
LaTeX and code artifacts, repeated compile-and-repair cycles, explicit theorem
and null-control demands, and reconstruction across conversations from handoff
packages. Those observations justify treating the collaboration as a serious
candidate case. They do not establish the complete certificate.

\begin{center}
\small
\begin{tabularx}{\textwidth}{@{}>{\raggedright\arraybackslash}p{0.22\textwidth}>{\raggedright\arraybackslash}p{0.18\textwidth}>{\raggedright\arraybackslash}X@{}}
\toprule
Gate family & Present evidential status & Conservative assessment \\
\midrule
Declared heterogeneous components & Observational & Human, model, artifacts,
and execution environment are distinguishable components with different
roles and information access. \\
Persistent shared vocabulary & Observational & Recurrent terms and distinctions
survive across research phases and constrain later work. \\
Exact structural core \(\Gamma(H)\) & Unmeasured & No frozen local coordinate
system or learned transfer family has yet been estimated from held-out data. \\
Transfer-Laplacian gap & Unmeasured & The idealized value \(2\) belongs to the
toy four-cycle, not to the actual collaboration. \\
Operational alignment & Unmeasured & No preregistered future variable has been
shown conditionally independent of the full history given a measured workspace
state. \\
Joint advantage & Suggestive only & The coupled process has produced outputs
that neither party appears to have supplied verbatim, but local-only Bayes risks
have not been measured under controlled information ceilings. \\
Matched relational null gap & Unmeasured & Shuffled artifacts, matched irrelevant
memory, and local-predictive-marginal-preserving nulls have not been run as a
frozen panel. \\
Causal workspace efficacy & Suggestive only & Corrections, compiler failures,
and artifact substitutions change later behavior, but they are not yet matched
workspace interventions with manipulation checks. \\
Artifact-mediated persistence & Observational & Handoff files and source bundles
support continuation, but blind reconstruction and ablation curves have not
been estimated. \\
Formation provenance & Mixed & The communication architecture is scaffolded,
while terminology, theorem organization, and many control practices were
adaptively developed. Causal formation-history contribution is unmeasured. \\
Workspace agency & Not established & No complete goal-retention,
counterfactual-policy, distributed-control, damage, and repair certificate has
been run. \\
Consciousness or personhood & Not established & The theory supplies no such
inference from the observed collaboration. \\
\bottomrule
\end{tabularx}
\end{center}

The correct present statement is therefore:
\[
\boxed{
\text{candidate persistent collaborative workspace}
\quad\text{rather than}\quad
\text{certified relational workspace agent}.
}
\]

\subtopic{79.6 A controlled certification design}

A serious test can use one manuscript section or software work unit as the unit of analysis.
Before the certification split is opened, freeze:
\begin{itemize}
\item the next-stage task family;
\item the downstream target, such as accepted theorem patch, successful clean
 compile, and absence of a declared regression;
\item the local information available to each component;
\item the workspace extractor and transfer maps;
\item all intervention and null families; and
\item the required gate thresholds.
\end{itemize}
The following conditions can then be compared:
\[
\begin{array}{ll}
\mathsf H & =\text{human-side externalized state only},\\
\mathsf M & =\text{fresh model instance with model-side state only},\\
\mathsf A & =\text{artifact package without live human or model state},\\
\mathsf X & =\text{executor and test outputs without the research context},\\
\mathsf Full & =\text{the intact recurrent loop}.
\end{array}
\]
The design must avoid pretending that a person's unobserved internal state has
been fully measured. The local conditions certify only the declared accessible
histories.

A minimal null panel includes:
\begin{enumerate}
\item theorem-to-artifact reassignment with artifact size and style matched;
\item temporal shuffling of verification feedback;
\item matched irrelevant manuals or source bundles;
\item context reset followed by blind reconstruction;
\item removal of compiler or test feedback while preserving message count;
\item model replacement under a fixed handoff package;
\item artifact removal with equivalent token budget restored by irrelevant text;
\item a local-marginal-matched synthetic null for the frozen downstream target.
\end{enumerate}
A valid empirical claim would require the intact loop to outperform every
required local and matched-null condition, while the extracted workspace state
remains sufficient and feasible workspace interventions alter the downstream
result in the predicted direction.

\subtopic{79.7 Handoff and identity experiment}

Let \(B_t\) be the source-and-results bundle at the end of stage \(t\). A blind
handoff trial replaces the active model instance and optionally the human
analyst, then asks the replacement coalition to reconstruct:
\[
\left(
\text{current definitions},
\text{proved results},
\text{open obligations},
\text{next admissible research step}
\right).
\]
The reconstruction error is measured against a frozen adjudication set. A
replacement invariance test additionally compares:
\begin{itemize}
\item the transported structural workspace;
\item operation behavior on held-out transformations;
\item response to matched interventions; and
\item preservation of failed and unresolved claims.
\end{itemize}
Equivalent benchmark prose alone is insufficient. A replacement that silently
changes which interventions alter later reasoning has a nonzero causal identity
defect even when it produces a similarly polished manuscript.

\subtopic{79.8 Case card}

For the idealized four-cycle model,
\[
\mathsf S,\mathsf O,\mathsf R,\mathsf P
\]
can all be made analytic by specifying the downstream kernel and temporal
updates. For the motivating collaboration itself, the conservative card is
\[
\boxed{
\begin{aligned}
\mathsf S&=\text{observational candidate},
&\mathsf O&=\text{unmeasured},\\
\mathsf R&=\text{unmeasured},
&\mathsf P&=\text{observational candidate},\\
\mathsf F&=\text{mixed and unmeasured},
&\mathsf A&=\text{not established},\\
\mathsf D&=\text{unmeasured},
&\mathsf C&=\text{not established}.
\end{aligned}
}
\]
This audit is intentionally stricter than ordinary conversational language. It
preserves the possibility that the collaboration instantiates a real relational
workspace while refusing to treat that possibility as already certified.

\topic{80}{Case Study C: A Closed-Form Multi-Agent Workspace Agent}

\subtopic{80.1 Declared finite-state system}

Consider three agents:
\[
V_{\mathrm{MA}}=\{P,M,R\},
\]
called the planner, monitor, and recorder. The global workspace-control state
is
\[
c_t=(g_t,b_t,h_t)\in\mathbb F_2^3,
\]
where:
\[
\begin{array}{ll}
g_t=\text{goal bit},
&b_t=\text{environment or belief bit},\\
h_t=\text{workspace-health bit},
&h_t=1\text{ means intact}.
\end{array}
\]
The local accessible states are
\[
X_P(t)=(g_t,b_t),
\qquad
X_M(t)=(b_t,h_t),
\qquad
X_R(t)=(g_t,h_t).
\]
This is the three-node distributed code from Topic~77. Consequently:
\[
\dim\Gamma=3,
\qquad
\spec(L)=\{0,0,0,2,2,2\},
\qquad
k_{\mathrm{rec}}=2,
\qquad
d_{\mathrm{er}}=2.
\]
The organization is fully scaffolded. No endogenous-formation claim is made.

\subtopic{80.2 Goal retention}

In the absence of an explicit goal intervention, let
\[
g_{t+1}=g_t.
\]
Under matched cue-cut trials,
\[
\operatorname{do}(g_t=0)
\quad\text{and}\quad
\operatorname{do}(g_t=1)
\]
remain distinguishable at every later time:
\[
|g_{t+\tau}-g_{t+\tau}'|=1.
\]
Thus the cue-controlled retention factor is exactly one:
\[
\eta_G(\tau)=1.
\]
The goal is not reloaded by an external channel after the initial intervention.

\subtopic{80.3 Joint task and repair action}

After a possible damage event, let the pre-repair health be
\[
h_t^-\in\{0,1\}.
\]
The system chooses a joint action
\[
u_t=(a_t,r_t)\in\{0,1\}^2,
\]
where \(a_t\) is a task action and \(r_t\) is a repair action. Define the
optimal task action
\[
a_t^*=g_t\oplus b_t\oplus h_t^-
\]
and optimal repair action
\[
r_t^*=1-h_t^-.
\]
No individual node determines the full pair \((a_t^*,r_t^*)\):
\begin{itemize}
\item the planner lacks \(h_t^-\);
\item the monitor lacks \(g_t\); and
\item the recorder lacks \(b_t\).
\end{itemize}
Every two-node coalition determines all three workspace bits and therefore the
complete action. Under independent uniform bits,
\[
H(a_t^*\mid X_i(t))=1
\qquad\text{for every }i,
\]
so the task-action component has a one-bit joint advantage and
\[
k_{\mathrm{ctl},0}=2.
\]

\subtopic{80.4 Interventional value and exact policy}

Let the one-step utility be
\[
U_g(a,r;b,h^-)
=
\mathbf 1\{a=g\oplus b\oplus h^-\}
+2\mathbf 1\{h^-\lor r=1\}
-\frac12r.
\]
For fixed \((g,b,h^-)\), the unique maximizer is
\[
\boxed{
(a^*,r^*)
=
(g\oplus b\oplus h^-,1-h^-).
}
\]
Indeed:
\begin{itemize}
\item a wrong task action loses one unit;
\item omitting repair when \(h^-=0\) loses \(3/2\) units; and
\item performing unnecessary repair when \(h^-=1\) loses \(1/2\) unit.
\end{itemize}
The smallest optimal-action gap is therefore
\[
\Delta_Q=\frac12.
\]
By Theorem~\ref{thm:counterfactual-policy}, any uniformly accurate value model with
\[
\varepsilon_Q<\frac14
\]
recovers the exact greedy joint action. More generally, its regret is bounded
by
\[
\operatorname{Reg}\leq2\varepsilon_Q+\delta_{\mathrm{opt}}.
\]

\subtopic{80.5 Feasible goal, workspace, and damage interventions}

The intervention family contains:
\[
\operatorname{do}(g_t=a),
\qquad
\operatorname{do}(b_t=b),
\qquad
\operatorname{do}(h_t^-=h),
\]
implemented consistently across every local chart that contains the manipulated
coordinate. Holding the other two bits fixed, each core-bit intervention
changes the prescribed task action. A matched intervention on an independent
residual bit leaves the task and repair actions unchanged. Hence
\[
\Delta_{\mathrm{specific}}>0.
\]
A damage intervention sets
\[
h_t^-=0.
\]
The matched no-repair counterfactual holds \(r_t=0\), while the intact policy
sets \(r_t=1\).

\subtopic{80.6 Exact self-maintenance}

Post-action health evolves according to
\[
h_{t+1}=h_t^-\lor r_t.
\]
Define the maintenance deficit
\[
e_t=1-h_t.
\]
Under the intact policy,
\[
h_{t+1}=1
\qquad\text{and}\qquad
e_{t+1}=0
\]
for both damage states. Under the matched no-repair intervention and
\(h_t^-=0\),
\[
e_{t+1}^{(0)}=1.
\]
Thus the measured causal repair gain is
\[
\boxed{
\Delta_{\mathrm{maint}}=1.
}
\]
The closed-loop recurrence has
\[
\beta=0,
\qquad
\overline d=0,
\qquad
\overline\epsilon=0,
\]
so Theorem~\ref{thm:self-maintenance} gives
\[
\sup_t e_t=0
\]
after each repair action. Monitor-blind random repair fails the specificity
control whenever it either omits necessary repair or performs unnecessary
repair with positive probability.

\subtopic{80.7 Reentry and causal loop}

The repaired health bit becomes part of the next workspace state:
\[
h_t^-
\longrightarrow
r_t
\longrightarrow
h_{t+1}
\longrightarrow
X_M(t+1),X_R(t+1).
\]
Cutting the direct storage of \(h_t^-\) while preserving the action-mediated
path still leaves a return path through repair. Cutting the repair-to-health
edge destroys that return. The system therefore contains explicit
loop-mediated reentry rather than mere temporal autocorrelation.

\subtopic{80.8 Analytic agency certificate}

Within this finite structural causal model, the following hold exactly:
\begin{enumerate}
\item a nontrivial exact distributed workspace with spectral gap two;
\item two-node minimal reconstruction and control coalition size;
\item matched-null joint advantage for the parity task;
\item persistent internally retained goal state;
\item exact interventional action values and zero policy regret;
\item goal and workspace interventions that alter action;
\item damage detection and repair selection;
\item positive matched repair gain and bounded maintenance deficit; and
\item causal reentry through the action-mediated repair loop.
\end{enumerate}
Therefore, relative to the declared bits, actions, utility, intervention family,
and one-step horizon,
\[
\boxed{
\operatorname{Agent}_{\Theta_A}(\Workspace)=1.
}
\]
The case reaches formal Level~6 by construction. It remains:
\[
\text{fully scaffolded},
\qquad
\text{domain restricted},
\qquad
\text{non-evidence for consciousness}.
\]

\subtopic{80.9 Case card}

\[
\boxed{
\begin{aligned}
\mathsf S&=\checkmark_{\mathrm{analytic}},
&\mathsf O&=\checkmark_{\mathrm{analytic}},\\
\mathsf R&=\checkmark_{\mathrm{analytic}},
&\mathsf P&=\checkmark_{\mathrm{analytic}},\\
\mathsf F&=\text{fully scaffolded},
&\mathsf A&=\checkmark_{\mathrm{analytic}},\\
\mathsf D&=\text{fixed decomposition},
&\mathsf C&=\text{not established}.
\end{aligned}
}
\]
The example demonstrates logical compatibility among relationality, redundancy,
agency, and self-maintenance. It does not show that those properties arise
spontaneously, scale to open environments, or imply subjective experience.

\topic{81}{Near-Miss Cases and Anti-Triviality Controls}

Positive examples are insufficient unless neighboring false positives are also
made explicit. The following systems each resemble a relational workspace in
one respect while failing a required gate.

\subtopic{81.1 Near miss A: the replicated archive}

Let every node store a complete copy of the same core:
\[
E_i=I_d
\qquad\text{for every }i.
\]
Then coherence is perfect and reconstruction is robust, but
\[
k_0=1,
\qquad
\min_i\delta_i^{\mathrm{local}}=0.
\]
If the archive is never used in downstream control, then also
\[
\eta_{\mathrm{causal}}=0.
\]
This is a shared representation or archive, not a relationally irreducible
workspace.

\subtopic{81.2 Near miss B: the centralized oracle with decorative agents}

Suppose node \(o\) contains the complete predictive state \(Z_t\), while every
other node receives only messages generated by \(o\). The full system can have
excellent performance, rich communication, and recurrent activity. Yet
\[
Y_t^+\indep\mathbf H_t\mid X_o(t),
\qquad
k_0=1,
\qquad
\Delta_{\mathrm{joint}}=0.
\]
Removing the decorative nodes may change latency or style without changing the
operative capacity. Multi-agent appearance therefore does not imply a
relational workspace.

\subtopic{81.3 Near miss C: locally coherent but globally obstructed charts}

Let three local scalar charts have relative overlap demands
\[
b_{12}=b_{23}=b_{31}=1.
\]
The cycle obstruction is
\[
\omega_{\mathcal U}(b)
=
\frac{|1+1+1|}{\sqrt3}
=\sqrt3>0.
\]
Each local chart is internally coherent, but no global scalar assignment
satisfies all three overlap equations. The system contains compatible-looking
subworkspaces without a global workspace.

\subtopic{81.4 Near miss D: externally reloaded goals}

Let
\[
G_{t+1}=C_{t+1},
\]
where \(C_{t+1}\) is an external instruction repeated at every step. Goal and
action can be strongly correlated, but after cutting the cue channel,
\[
I(G_{t+\tau};G_t\mid\operatorname{cut}(C))=0.
\]
The system tracks an external command. It has not established internal goal
retention.

\subtopic{81.5 Near miss E: task recovery with identity failure}

Let two realizations have identical benchmark output distributions,
\[
P(Y_{\mathrm{bench}})=P'(Y_{\mathrm{bench}}),
\]
but different intervention kernels:
\[
K(w,b)\neq K'(\beta w,b)
\]
for some matched workspace intervention. Then task behavior is preserved while
\[
d_{\mathrm{causal}}(W,W')>0.
\]
The later system may be a successor or replacement rather than the same
workspace under the identity criterion of Part~XII.

\subtopic{81.6 Near miss F: open-loop repair replay}

A recorded repair sequence may restore average task performance after a fixed
damage schedule. If the same sequence is replayed when the damage timing or
type changes, it need not respond appropriately. When
\[
R_t\indep M_t\mid\text{clock and replay index},
\]
monitor-to-repair mediation is absent. Successful replay is robustness under a
script, not self-maintenance.

\subtopic{81.7 Anti-triviality summary}

The near misses establish six independent exclusions:
\[
\boxed{
\begin{aligned}
\text{workspace}
\neq{}&\text{replicated storage},\\
\neq{}&\text{centralized control with decorative communication},\\
\neq{}&\text{locally coherent but obstructed charts},\\
\neq{}&\text{externally reloaded goals},\\
\neq{}&\text{task-equivalent identity replacement},\\
\neq{}&\text{open-loop repair replay}.
\end{aligned}
}
\]
A complete certificate must distinguish the intended realization from all
near misses relevant to its claimed level.

\topic{82}{Cross-Case Synthesis and Benchmark Blueprint}

\subtopic{82.1 Comparative gate matrix}

\begin{center}
\small
\begin{tabularx}{\textwidth}{@{}>{\raggedright\arraybackslash}p{0.20\textwidth}>{\raggedright\arraybackslash}p{0.20\textwidth}>{\raggedright\arraybackslash}p{0.23\textwidth}>{\raggedright\arraybackslash}X@{}}
\toprule
Gate family & Human--human analytic model & Motivating human--AI collaboration & Multi-agent analytic model \\
\midrule
Exact structural core & Analytic pass & Not measured & Analytic pass \\
Spectral gap & \(2\) & Not measured & \(2\) \\
Minimal sufficient coalition & \(k_0=2\) & Not measured & \(k_0=2\) \\
Matched-null transfer gap & \(1\) bit & Not measured & \(1\) bit for task parity \\
Causal specificity & Analytic pass & Suggestive, uncontrolled & Analytic pass \\
One-node erasure resilience & Analytic pass & Handoff only observational & Analytic pass \\
Formation provenance & Scaffolded & Mixed, unmeasured & Fully scaffolded \\
Goal retention & Not modeled & Not certified & Exact analytic pass \\
Counterfactual policy & Not modeled & Not certified & Exact analytic pass \\
Self-maintenance & Not modeled & Not certified & Exact analytic pass \\
Highest supported level & Level 4 in model & Candidate only & Level 6 in model \\
Consciousness/\allowbreak personhood & Not established & Not established & Not established \\
\bottomrule
\end{tabularx}
\end{center}

The table illustrates why the theory uses a vector certificate rather than one
score. The human--human model is relational and persistent without being
agentive. The multi-agent model is agentive without being endogenously formed.
The motivating collaboration is richer than either toy model but presently has
weaker evidential status because the necessary interventions and nulls have not
been frozen and run.

\subtopic{82.2 Benchmark family B1: distributed parity}

The first benchmark is the exact parity family from Topics~85--80. For
\(m\geq3\), distribute \(m\) independent bits among nodes so that no individual
node sees all bits, but specified coalitions do. The target is
\[
Y=\bigoplus_{j=1}^m C_j.
\]
The benchmark measures:
\begin{itemize}
\item the learned transfer-Laplacian spectrum;
\item coalition conditional entropy and log risk;
\item recovery of \(k_0\);
\item robustness to edge and node ablations;
\item matched local-predictive null gaps; and
\item causal response to consistent versus inconsistent coordinate
 interventions.
\end{itemize}
The ground truth is known exactly, so structural and relational estimators can
be calibrated before they are applied to open-ended research tasks.

\subtopic{82.3 Benchmark family B2: theorem-patch handoff}

A theorem-patch task supplies:
\[
\left(
\text{definitions},
\text{partial proof},
\text{counterexample tests},
\text{style constraints}
\right)
\]
across distinct components. No component receives all four. The coupled
system must produce a patch that:
\begin{enumerate}
\item is formally consistent with the definitions;
\item closes the declared proof obligation or marks it unresolved;
\item passes the frozen counterexample suite;
\item compiles from a clean environment; and
\item survives a blind replacement-handoff trial.
\end{enumerate}
Matched nulls permute proof fragments, substitute irrelevant artifacts of equal
size, or replay test outputs from a different theorem. This family targets the
structural, relational, persistence, and identity layers without requiring an
agency claim.

\subtopic{82.4 Benchmark family B3: damage and repair}

A certified workspace is first formed and measured. A damage intervention then
removes or corrupts one of:
\[
\text{a transfer channel},
\quad
\text{a participant},
\quad
\text{an artifact},
\quad
\text{a goal encoding},
\quad
\text{a verification path}.
\]
The system receives a repair-action family with matched costs. The benchmark
measures:
\[
\eta_{D\to M},
\qquad
\eta_{M\to R},
\qquad
\eta_{R\to v},
\qquad
\Delta_{\mathrm{repair}}^{\mathrm{spec}},
\qquad
\mathbf d_{\mathrm{id}}.
\]
Open-loop replay, random repair, off-workspace repair, and identity-breaking
task recovery are mandatory controls.

\subtopic{82.5 Benchmark family B4: scaffold transport}

Train the same formation rule under several scaffolds
\[
\Sigma^{(1)},\ldots,\Sigma^{(K)}
\]
that differ in node labels, coordinate systems, and communication topology but
support the same abstract task. After gauge alignment, measure whether the
learned organizations transport to a common equivalence class:
\[
[\widehat H^{(1)}]_{\mathcal G}
\cong\cdots\cong
[\widehat H^{(K)}]_{\mathcal G}.
\]
A structure that appears only under one answer-encoding scaffold has weak
endogeneity evidence. A structure that recurs across alternative scaffolds,
passes held-out task and causal gates, and cannot be compressed into the
scaffold specification has stronger evidence.

\subtopic{82.6 Reporting rule for mixed analytic and empirical cases}

When a case combines exact mathematics with observed data, report separate
columns:
\[
\boxed{
\text{proved in model}
\quad\big|\quad
\text{estimated in system}
\quad\big|\quad
\text{not established}.
}
\]
For example, the four-cycle Laplacian has an analytic gap of two, but an actual
human--AI collaboration has that gap only if its measured discrepancy operator
has the corresponding spectrum. Importing the toy value into the empirical
case is invalid.

\subtopic{82.7 Lessons from the integrated cases}

The case studies establish the following conceptual separation:
\begin{enumerate}
\item A relational workspace can be mathematically exact without being an
 agent.
\item An agentive workspace can be fully scaffolded and still be operationally
 real within its domain.
\item A rich real collaboration can be a stronger candidate than a toy model
 while having a weaker current certificate.
\item Redundancy and irreducibility are compatible when no individual node is
 sufficient but several proper coalitions are.
\item Artifact-mediated persistence is not passive storage when artifact
 contents causally alter later operations.
\item High task performance, communication volume, or behavioral polish alone
 cannot distinguish the intended theory from its near misses.
\item No structural, causal, or agentive result in these examples entails
 phenomenal consciousness.
\end{enumerate}

\centerbox{\large\bfseries
The theory becomes empirically meaningful only when the same gate card that
certifies the positive case rejects its matched near misses.
}
